\section{Decision Framework}
\label{sec:decision}

\subsection{The Four-Question Tree}

The comparison in Sections~\ref{sec:compactness}
and~\ref{sec:distributional} reveals that practitioners can select an
appropriate algorithm by answering four binary questions in sequence.
The decision tree below maps these questions to algorithm recommendations.
The tree encodes structural properties of the algorithms (established in
G.6--G.13), not solely the single-run empirical results of this paper.

\begin{enumerate}[label=\textbf{Q\arabic*.}]

  \item \textbf{Is VRA compliance required?}

    If the redistricting task involves states or jurisdictions subject to
    Section~5 or Section~2 of the Voting Rights Act, and minority-majority
    district preservation must be enforced during sampling (not merely
    verified post hoc), then VRA-Aware ReCom (G.13) is the only algorithm
    that guarantees preservation of minority-majority districts
    throughout the chain via adaptive boosts.
    \emph{Proceed to: \vra{} (G.13). Stop.}

    If VRA compliance is verified post hoc (acceptable for many contexts),
    proceed to Q2.

  \item \textbf{Is $k \geq 20$?}

    For states with $k \geq 20$ congressional districts (TX $k=38$,
    FL $k=28$, CA $k=52$, NY $k=26$, PA $k=17$, IL $k=17$, OH $k=15$,
    GA $k=14$ borderline, and any state approaching this threshold),
    single-scale chains mix slowly.
    The lag-100 autocorrelation results in Section~\ref{sec:distributional}
    show that Multi-scale (G.11) provides a meaningful mixing improvement
    for TX; for CA ($k=52$), the improvement is expected to be larger.
    \emph{Proceed to: \ms{} (G.11). Stop.}

    For $k < 20$, single-scale chains mix adequately; proceed to Q3.

  \item \textbf{Is a calibrated target distribution required?}

    If the ensemble will be used for litigation support, academic
    publication under peer review, or any context where the target
    distribution must be mathematically specified, a calibrated distribution
    is required.

    \begin{itemize}
      \item For \emph{exact} calibration (importance-weighted particles,
            user-specified target): use \smc{} (G.7).
      \item For \emph{approximate} calibration (approximately uniform over
            valid plans, Metropolis-Hastings corrected): use \fr{} (G.9)
            or \msp{} (G.10).
            \msp{} is recommended as the default: it has marginally lower
            autocorrelation and similar distributional guarantees.
    \end{itemize}

    If distributional calibration is not required (e.g., the task is
    compactness optimisation or sensitivity analysis), proceed to Q4.

  \item \textbf{Is pure compactness optimisation the goal?}

    If the goal is to produce the most compact plan achievable within
    the population-balance constraint (for example, constructing a
    benchmark for expert analysis or a starting point for further
    refinement):

    \begin{itemize}
      \item For fastest compactness optimisation with no distributional
            constraint: use \sburst{} (G.6).
      \item For compactness optimisation with a better-characterised
            distributional footprint (MH-corrected): use
            \sburstf{} (G.12).
    \end{itemize}

    If the goal is not compactness optimisation, the remaining algorithm
    is \flip{} (G.8): the appropriate choice for local boundary sensitivity
    analysis at a fixed initial plan.
    \emph{Proceed to: \flip{} (G.8). Stop.}

\end{enumerate}

\subsection{Decision Tree Diagram}

\begin{center}
\begin{tikzpicture}[
  decision/.style={
    diamond, draw, fill=blue!8, text width=3.8cm, align=center,
    inner sep=2pt, font=\small
  },
  result/.style={
    rectangle, draw, fill=green!10, text width=3.2cm, align=center,
    rounded corners=3pt, font=\small
  },
  arrow/.style={->, >=Stealth, thick},
  label/.style={font=\small\itshape},
]

%% Root
\node[decision] (q1) {VRA compliance\\required?};

%% Q1 Yes
\node[result, right=2.0cm of q1, yshift=1.4cm] (vra)
  {\vra{} (G.13)\\\footnotesize VRA-constrained};
\draw[arrow] (q1.north east) -- node[label, above, sloped]{\small Yes} (vra.west);

%% Q2
\node[decision, below=1.6cm of q1] (q2) {$k \geq 20$?};
\draw[arrow] (q1.south) -- node[label, right]{\small No} (q2.north);

%% Q2 Yes
\node[result, right=2.0cm of q2, yshift=1.0cm] (ms)
  {\ms{} (G.11)\\\footnotesize Large-$k$ mixing};
\draw[arrow] (q2.north east) -- node[label, above, sloped]{\small Yes} (ms.west);

%% Q3
\node[decision, below=1.6cm of q2] (q3) {Calibrated\\distribution\\required?};
\draw[arrow] (q2.south) -- node[label, right]{\small No} (q3.north);

%% Q3 Yes — two outcomes: Exact (SMC) or Approximate (MergeSplit/ForestReCom)
\node[result, right=2.0cm of q3, yshift=1.4cm] (smc)
  {\smc{} (G.7)\\\footnotesize Exact calibration};
\node[result, right=2.0cm of q3, yshift=-0.5cm] (ms2)
  {\msp{}/\fr{} (G.10/9)\\\footnotesize Approx.\ uniform};
%% Branch point east of q3
\coordinate (branch3) at ($(q3.east) + (0.8cm, 0)$);
\draw[arrow] (q3.east) -- node[label, above]{\small Yes} (branch3);
\draw[arrow] (branch3) |- node[label, right, pos=0.75]{\small Exact} (smc.west);
\draw[arrow] (branch3) |- node[label, right, pos=0.75]{\small Approx} (ms2.west);

%% Q4
\node[decision, below=1.6cm of q3] (q4) {Compactness\\optimisation?};
\draw[arrow] (q3.south) -- node[label, right]{\small No} (q4.north);

%% Q4 Yes
\node[result, right=2.0cm of q4, yshift=1.2cm] (sb)
  {\sburst{} / \sburstf{} (G.6/12)\\\footnotesize Compact plans};
\draw[arrow] (q4.north east) -- node[label, above, sloped]{\small Yes} (sb.west);

%% Q4 No -> Flip
\node[result, below=1.0cm of q4] (flip)
  {\flip{} (G.8)\\\footnotesize Sensitivity analysis};
\draw[arrow] (q4.south) -- node[label, right]{\small No} (flip.north);

\end{tikzpicture}
\end{center}

\subsection{Summary Decision Matrix}

Table~\ref{tab:decision} distils the decision tree into a use-case matrix:
one recommended algorithm per use case with an alternative where applicable.
The alternative is appropriate when the primary recommendation is unavailable
(e.g., SMC particle budget is prohibitive) or when the practitioner has
domain-specific reasons to prefer a different distributional profile.

\begin{table}[ht]
\centering
\caption{%
  Decision matrix: recommended algorithm by use case.
  ``Alternative'' is the second choice when the primary recommendation
  is unavailable or inappropriate for the specific context.
}
\label{tab:decision}
\begin{tabular}{@{}lll@{}}
\toprule
Use Case & Recommended & Alternative \\
\midrule
Court submission (uniform distribution)
  & \smc{} (G.7) & \fr{} (G.9) \\
Compactness optimisation
  & \sburst{} (G.6) & \sburstf{} (G.12) \\
Large-$k$ states (TX, CA, FL)
  & \ms{} (G.11) & \sburst{} (G.6) \\
VRA-compliant sampling
  & \vra{} (G.13) & --- \\
Local sensitivity / boundary analysis
  & \flip{} (G.8) & --- \\
General purpose (small-$k$, no special requirements)
  & \msp{} (G.10) & \fr{} (G.9) \\
Compactness + distributional awareness
  & \sburstf{} (G.12) & \msp{} (G.10) \\
Academic publication (calibrated)
  & \smc{} (G.7) & \msp{} (G.10) \\
\bottomrule
\end{tabular}
\end{table}

\subsection{Three-Layer Compositor Integration}

The decision tree selects the \texttt{--search} algorithm.
In practice, algorithm selection interacts with the other two compositor
layers:

\paragraph{Structure (\texttt{--structure}).}
The initial plan produced by the structure layer is the starting point for
all chain-based methods.
A better initial plan (e.g., from \texttt{ratio-optimal} GeoSection
bisection) reduces the warm-up period for all chain methods.
For \smc{}, the structure layer is not used (SMC constructs plans from
scratch via particle filter).

\paragraph{Weights (\texttt{--weights-override}).}
Edge weights affect the probability of each spanning-tree cut in \recom{},
Forest ReCom, Merge-Split, and Multi-scale.
County weights (\texttt{county}) tend to produce plans that respect county
boundaries, regardless of the search algorithm.
VRA weights (\texttt{vra-aligned}) bias cuts toward minority community
boundaries; these are complementary to, but distinct from, the VRA-Aware
search mode.

\paragraph{Combining layers.}
The most powerful configurations combine all three layers.
For example:
\begin{itemize}
  \item \texttt{--structure prime-factor --weights county --search multiscale}:
        ApportionRegions initial plan + county-respecting weights +
        Multi-scale mixing.
        Recommended for TX or CA litigation.
  \item \texttt{--structure ratio-optimal --weights geographic --search vra-recom}:
        GeoSection initial plan + neutral weights + VRA-constrained chain.
        Recommended for VRA Section~2 analysis.
  \item \texttt{--structure ratio-optimal --weights geographic --search multi}:
        GeoSection initial plan + neutral weights + Short-Burst optimisation.
        Recommended for producing the most compact feasible plan.
\end{itemize}

\subsection{Limitations of the Decision Tree}

The decision tree presented here is based on the algorithms as implemented
in \texttt{redist-ensemble} as of G-series papers G.6--G.13.
Several caveats apply.

\textbf{Parameter sensitivity.}
Each algorithm has internal parameters (burst length for Short-Burst,
$\alpha$ for Multi-scale, VRA boost strength for VRA-Aware) that affect
performance.
The decision tree assumes default parameters; practitioners with unusual
geometries or strict distributional requirements should consult the
individual algorithm papers (G.6--G.13) for parameter guidance.

\textbf{State-specific effects.}
The comparison covers NC, WI, and TX.
States with unusual geographic features (e.g., California with very
heterogeneous tract sizes, or Alaska with non-contiguous population
centers) may require adjustment.

\textbf{Metric weights.}
The decision tree treats compactness optimisation and distributional
correctness as separate goals.
In practice, a practitioner may want some compactness improvement with
some distributional coverage.
ShortBurstForest (G.12) is designed for exactly this middle ground, but
the optimal balance point depends on the specific task.

\textbf{Future algorithms.}
CVD (B.22) and BFS Growth (B.23) will add geometric algorithms to the
comparison.
When these are implemented, the decision tree will be updated accordingly.
